% ============================================================
% CHAPITRE 1 — Introduction au design graphique
% ============================================================

\chapter{Introduction au design graphique}

\noindent
\textit{Ce premier chapitre pose les fondations de ta formation. Avant de manipuler le moindre outil, avant d'ouvrir Photoshop ou Illustrator, avant même de penser à créer une affiche ou un logo, il est indispensable de comprendre ce qu'est réellement le design graphique, d'où il vient, à quoi il sert, comment il se pratique. Ce chapitre te donnera les repères essentiels qui te serviront de boussole tout au long de ton apprentissage.}

\vspace{1em}

\noindent
Nous commencerons par définir le design graphique, puis nous retracerons son histoire, nous examinerons son rôle dans la communication contemporaine, nous cartographierons ses différents domaines, nous préciserons les missions et responsabilités du graphiste, nous distinguerons le design graphique des disciplines voisines, nous décrirons les qualités et compétences nécessaires, et nous détaillerons le processus de création, de l'idée au produit final.

% ============================================================
\section{Définition du design graphique}
% ============================================================

\subsection{Pourquoi commencer par une définition ?}

Avant de manipuler le moindre outil, avant d'ouvrir Photoshop ou Illustrator, avant même de penser à créer une affiche ou un logo, il est indispensable de comprendre ce que l'on fait réellement. Beaucoup de débutants confondent le design graphique avec le dessin, la décoration, l'art ou la simple maîtrise technique de logiciels. Cette confusion est la source de nombreuses erreurs qui se retrouvent ensuite dans les portfolios : des visuels certes « jolis », mais qui ne communiquent rien, qui ne répondent à aucun besoin, qui ne respectent aucune contrainte.

Cette première section a donc un objectif précis : poser une définition claire, complète et opérationnelle du design graphique. Non pas une définition vague et poétique, mais une définition qui te servira de boussole tout au long de ta formation et, plus tard, dans ta pratique professionnelle.

Nous allons procéder par étapes : d'abord une définition générale, puis une décomposition de ses termes, ensuite une définition professionnelle plus précise, et enfin une série de distinctions essentielles qui t'aideront à situer le design graphique par rapport à des disciplines voisines.

\subsection{Définition générale}

\begin{quote}
\textbf{Le design graphique est l'activité consistant à concevoir, organiser et produire des messages visuels destinés à transmettre une information, une idée ou une émotion à un public déterminé, en utilisant des éléments graphiques tels que la typographie, les images, les formes, les couleurs et la mise en page.}
\end{quote}

Cette définition, en apparence simple, contient en réalité plusieurs notions fondamentales qu'il faut examiner une par une. Reprenons-la terme par terme.

\subsubsection{« Concevoir »}

Le design graphique n'est pas d'abord un acte de production, c'est un acte de \terme{conception}. Concevoir signifie :

\begin{itemize}
    \item \terme{Penser} avant d'agir : réfléchir au problème à résoudre, au public à toucher, au message à faire passer.
    \item \terme{Planifier} : déterminer une stratégie visuelle, choisir des moyens adaptés, anticiper le résultat.
    \item \terme{Structurer} : organiser les éléments dans un ensemble cohérent, hiérarchisé, lisible.
\end{itemize}

C'est une distinction capitale. Un graphiste n'est pas quelqu'un qui « fait des images ». C'est quelqu'un qui \terme{résout des problèmes de communication par des moyens visuels}. Le logiciel n'est qu'un outil, un moyen. La conception, elle, se fait d'abord dans la tête, sur le papier, dans une réflexion stratégique.

Prenons un exemple concret. Une entreprise te demande de créer une affiche pour promouvoir un concert de jazz. Deux attitudes sont possibles :

\begin{itemize}
    \item \terme{L'attitude du producteur d'images} : ouvrir Photoshop, chercher une belle photo de saxophone, écrire le nom du groupe en gros, ajouter quelques effets, exporter. Résultat : une affiche peut-être agréable, mais générique, sans réflexion sur le public, sans hiérarchie claire, sans identité.
    \item \terme{L'attitude du designer graphique} : se demander qui est le public (amateurs de jazz ? étudiants ? familles ?), quel est le message principal (la date ? l'ambiance ? le nom de l'artiste ?), quel ton adopter (élégant, festif, intimiste ?), quels supports seront utilisés (affichage urbain ? réseaux sociaux ? programme imprimé ?). Ensuite seulement, concevoir une composition qui répond à ces questions.
\end{itemize}

La différence entre les deux n'est pas une différence de talent, c'est une différence de \terme{méthode}. C'est cette méthode que tu vas apprendre.

\subsubsection{« Organiser »}

Le design graphique est un art de l'\terme{organisation}. Une composition graphique n'est pas un empilement d'éléments posés au hasard : c'est un ensemble structuré où chaque élément a une place, une fonction, une raison d'être.

Organiser, cela signifie :

\begin{itemize}
    \item \terme{Hiérarchiser} : tout n'a pas la même importance. Il faut distinguer le message principal, les messages secondaires, les informations pratiques.
    \item \terme{Aligner} : les éléments doivent entretenir des relations visuelles cohérentes entre eux.
    \item \terme{Rythmer} : créer des respirations, des répétitions, des variations qui guident l'œil du lecteur.
    \item \terme{Équilibrer} : répartir les masses visuelles de manière harmonieuse, sans surcharger ni vider excessivement.
\end{itemize}

Cette notion d'organisation est ce qui distingue un visuel professionnel d'un visuel amateur. Un amateur va « remplir » un espace. Un professionnel va « construire » un espace.

\subsubsection{« Produire »}

Le design graphique aboutit toujours à une \terme{production}. C'est une discipline appliquée, pas une discipline purement théorique. Le résultat prend la forme d'un objet concret :

\begin{itemize}
    \item Un fichier numérique (image, PDF, fichier vectoriel).
    \item Un objet imprimé (affiche, brochure, carte de visite, packaging).
    \item Un support numérique (bannière web, publication pour réseaux sociaux, interface).
\end{itemize}

Cette dimension de production implique des contraintes techniques : résolution, format, colorimétrie, procédé d'impression, compatibilité des fichiers. Un bon designer graphique doit connaître ces contraintes, car une création magnifique mais techniquement inexploitable est un échec professionnel.

\subsubsection{« Des messages visuels »}

Le design graphique est un \terme{langage visuel}. Il ne s'agit pas de décorer, mais de \terme{communiquer}. Un message visuel peut avoir plusieurs fonctions :

\begin{itemize}
    \item \terme{Informer} : transmettre une information factuelle (une date, un lieu, un prix, un mode d'emploi).
    \item \terme{Identifier} : rendre une marque, un produit ou une organisation reconnaissable (logo, charte graphique).
    \item \terme{Promouvoir} : inciter à une action (acheter, s'inscrire, venir, voter).
    \item \terme{Émouvoir} : susciter une réaction affective (rire, surprise, nostalgie, admiration).
    \item \terme{Expliquer} : rendre compréhensible une notion complexe (infographie, schéma, signalétique).
\end{itemize}

Un même visuel peut cumuler plusieurs de ces fonctions, mais il doit toujours y avoir une \terme{intention communicative} claire. Si tu ne peux pas expliquer en une phrase ce que ton visuel cherche à dire ou à provoquer, c'est probablement qu'il ne communique rien.

\subsubsection{« À un public déterminé »}

C'est l'une des notions les plus importantes et l'une des plus négligées par les débutants. Le design graphique ne s'adresse jamais à « tout le monde » en général. Il s'adresse à un \terme{public précis}, avec ses habitudes, ses attentes, ses codes, sa culture visuelle.

Concevoir pour un public, cela implique :

\begin{itemize}
    \item \terme{Connaître ce public} : âge, milieu social, niveau culturel, habitudes de consommation des médias.
    \item \terme{Adapter le ton} : un visuel destiné à des enfants ne ressemblera pas à un visuel destiné à des cadres bancaires.
    \item \terme{Respecter les codes} : chaque secteur a ses conventions visuelles. Les ignorer, c'est prendre le risque de ne pas être compris.
    \item \terme{Se distinguer sans exclure} : le design doit être original, mais pas au point de perdre le public visé.
\end{itemize}

Un design qui ne s'adresse à personne ne touche personne. C'est une règle fondamentale.

\subsubsection{« En utilisant des éléments graphiques »}

Le design graphique dispose d'un vocabulaire visuel spécifique, composé de plusieurs familles d'éléments :

\begin{itemize}
    \item \terme{La typographie} : les caractères, les mots, les textes. C'est l'élément premier du design graphique, car il porte le message verbal.
    \item \terme{Les images} : photographies, illustrations, pictogrammes, dessins, textures.
    \item \terme{Les formes} : rectangles, cercles, lignes, tracés, aplats de couleur.
    \item \terme{Les couleurs} : teintes, valeurs, saturations, harmonies.
    \item \terme{La mise en page} : organisation spatiale de tous ces éléments dans un format donné.
\end{itemize}

Ces éléments ne sont pas indépendants : ils forment un \terme{système}. Le travail du designer consiste à les combiner de manière cohérente pour produire du sens.

\subsection{Définition professionnelle}

La définition générale que nous venons d'analyser peut être complétée par une définition plus opérationnelle, centrée sur l'activité professionnelle :

\begin{quote}
\textbf{Le design graphique est une discipline de communication visuelle qui consiste à résoudre un problème de transmission d'information ou d'émotion, en concevant des solutions visuelles adaptées à un public, à un support et à des contraintes techniques, dans le but de produire un effet déterminé.}
\end{quote}

Cette définition met l'accent sur trois notions essentielles :

\begin{enumerate}
    \item \terme{Le design graphique est une réponse à un problème.} Il n'existe pas « en soi ». Il naît d'un besoin : faire connaître un produit, expliquer un service, identifier une entreprise, célébrer un événement.
    \item \terme{Le design graphique est soumis à des contraintes.} Contraintes de format, de budget, de délai, de support, de droit d'auteur, de technologie. Ces contraintes ne sont pas des obstacles : elles font partie du travail de conception.
    \item \terme{Le design graphique vise un effet.} Il ne s'agit pas seulement de « faire beau ». Il s'agit de produire un résultat mesurable ou perceptible : attirer l'attention, faire comprendre, faire agir, faire mémoriser.
\end{enumerate}

\subsection{Ce que le design graphique n'est pas}

Pour bien comprendre ce qu'est le design graphique, il est utile de préciser ce qu'il n'est pas. Ces distinctions sont essentielles pour éviter les confusions les plus courantes.

\subsubsection{Le design graphique n'est pas de l'art}

L'art et le design graphique partagent des moyens (formes, couleurs, compositions), mais leurs finalités sont différentes.

\begin{table}[h!]
\centering
\begin{tabular}{@{}p{6cm}p{6cm}@{}}
\toprule
\textbf{Art} & \textbf{Design graphique} \\
\midrule
Expression personnelle de l'artiste & Réponse à un besoin de communication \\
L'artiste choisit son sujet & Le designer reçoit un brief \\
L'œuvre peut être ambiguë, ouverte & Le message doit être compris \\
Le public est celui qui veut bien regarder & Le public est ciblé et défini \\
Pas de contrainte de délai ou de format (en général) & Contraintes techniques, économiques, temporelles \\
La valeur est esthétique et symbolique & La valeur est fonctionnelle et communicative \\
\bottomrule
\end{tabular}
\caption{Comparaison entre art et design graphique.}
\end{table}

Cela ne signifie pas que le design graphique ne peut pas être beau, ni qu'il ne peut pas être artistique. Cela signifie que sa \terme{finalité première n'est pas l'expression de soi}, mais la \terme{communication}. Un designer qui crée pour se faire plaisir, sans tenir compte du public et du message, fait de l'art (ce qui est respectable), mais pas du design graphique.

\subsubsection{Le design graphique n'est pas du dessin}

Savoir dessiner est une compétence utile, mais ce n'est pas une condition nécessaire pour être designer graphique. Beaucoup de designers professionnels dessinent peu ou mal. En revanche, ils savent :

\begin{itemize}
    \item Structurer une composition.
    \item Choisir et associer des typographies.
    \item Hiérarchiser l'information.
    \item Utiliser la couleur à bon escient.
    \item Comprendre un brief et y répondre.
\end{itemize}

Le dessin est un outil parmi d'autres, pas la compétence centrale. Un designer peut travailler uniquement avec de la typographie, des formes géométriques et des photographies, sans jamais dessiner.

\subsubsection{Le design graphique n'est pas de la décoration}

Décorer consiste à ajouter des éléments pour embellir un support existant. Le design graphique, lui, \terme{construit} le support dans son ensemble, en fonction d'un objectif. La décoration est accessoire ; le design est structurel.

Un designer ne se demande pas « qu'est-ce que je pourrais ajouter pour que ce soit plus joli ? », mais « qu'est-ce qui est nécessaire pour que le message passe clairement et efficacement ? ».

\subsubsection{Le design graphique n'est pas de la simple maîtrise logicielle}

Savoir utiliser Photoshop et Illustrator est indispensable, mais ce n'est pas suffisant. Un technicien qui maîtrise les logiciels mais ne comprend pas les principes de composition, de typographie, de couleur et de communication ne produira que des visuels vides. La maîtrise technique est au service de la conception, pas l'inverse.

\subsection{Les quatre piliers du design graphique}

Pour synthétiser, on peut dire que le design graphique repose sur quatre piliers fondamentaux.

\subsubsection{Pilier 1 : La fonction}

Tout design graphique a une \terme{fonction} : informer, identifier, promouvoir, expliquer, émouvoir. Cette fonction détermine les choix graphiques. Un design qui ne remplit pas sa fonction, même s'il est esthétiquement réussi, est un échec.

\textbf{Question à se poser :} à quoi sert ce visuel ? Que doit-il accomplir ?

\subsubsection{Pilier 2 : La forme}

La \terme{forme} est l'ensemble des moyens visuels mis en œuvre : composition, typographie, couleur, image, format. C'est la dimension esthétique et technique du design. Elle doit être au service de la fonction, et non l'inverse.

\textbf{Question à se poser :} comment organiser les éléments pour atteindre l'objectif ?

\subsubsection{Pilier 3 : Le public}

Le \terme{public} est la destination du message. Le design graphique n'existe pas dans l'absolu : il est toujours reçu par quelqu'un. Connaître ce quelqu'un est essentiel pour adapter le ton, le niveau de langage, les codes visuels.

\textbf{Question à se poser :} à qui s'adresse ce visuel ? Que sait-il déjà ? Qu'attend-il ?

\subsubsection{Pilier 4 : Le contexte}

Le \terme{contexte} englobe les contraintes techniques, économiques, culturelles et temporelles dans lesquelles le design est produit. Un design réussi dans un contexte peut échouer dans un autre. Le contexte inclut : le support (papier, écran, objet), le format, le budget, le délai, la culture du pays ou du secteur.

\textbf{Question à se poser :} dans quelles conditions ce visuel sera-t-il produit, diffusé et vu ?

\subsection{Exemples concrets d'application de la définition}

Pour rendre cette définition plus concrète, examinons trois exemples de projets graphiques et voyons comment les quatre piliers s'y appliquent.

\subsubsection{Exemple 1 : Un logo d'entreprise}

\begin{itemize}
    \item \terme{Fonction} : identifier l'entreprise, la rendre reconnaissable, transmettre ses valeurs (confiance, modernité, tradition, etc.).
    \item \terme{Forme} : un symbole simple, une typographie choisie, une couleur dominante, un agencement équilibré.
    \item \terme{Public} : les clients, les partenaires, les employés. Le logo doit être compris et mémorisé par eux.
    \item \terme{Contexte} : le logo sera utilisé sur des supports très variés (papier, écran, tissu, gravure). Il doit donc être lisible à toutes les tailles, en couleur et en noir et blanc, et fonctionner dans différents contextes.
\end{itemize}

\subsubsection{Exemple 2 : Une affiche de film}

\begin{itemize}
    \item \terme{Fonction} : attirer l'attention, donner envie de voir le film, communiquer le genre (action, drame, comédie), informer sur la date de sortie.
    \item \terme{Forme} : une image forte (souvent le visage de l'acteur principal), un titre lisible, une ambiance colorée, une composition dramatique.
    \item \terme{Public} : les spectateurs potentiels. L'affiche doit parler aux amateurs du genre, mais aussi séduire un public plus large.
    \item \terme{Contexte} : affichage urbain (vu de loin, en quelques secondes), mais aussi reproduction en petit format (presse, web). L'affiche doit fonctionner à plusieurs échelles.
\end{itemize}

\subsubsection{Exemple 3 : Une infographie explicative}

\begin{itemize}
    \item \terme{Fonction} : rendre compréhensible une information complexe (statistiques, processus, comparaison).
    \item \terme{Forme} : schémas, pictogrammes, graphiques, couleurs distinctives, hiérarchie claire.
    \item \terme{Public} : lecteurs non spécialistes. L'infographie doit être compréhensible sans connaissance préalable du sujet.
    \item \terme{Contexte} : publication dans un journal, un rapport, un site web. Le format et la résolution doivent être adaptés.
\end{itemize}

Dans chacun de ces exemples, on retrouve les mêmes questions : à quoi ça sert ? pour qui ? comment ? dans quelles conditions ? C'est cela, la démarche du design graphique.

\subsection{Le vocabulaire essentiel de cette section}

Pour fixer les notions, voici les termes clés à retenir et à comprendre. Tu les retrouveras tout au long de ta formation.

\begin{table}[h!]
\centering
\begin{tabular}{@{}p{4cm}p{10cm}@{}}
\toprule
\textbf{Terme} & \textbf{Définition} \\
\midrule
Design graphique & Discipline de conception de messages visuels destinés à communiquer une information, une idée ou une émotion à un public déterminé. \\
Communication visuelle & Ensemble des moyens visuels utilisés pour transmettre un message (image, typographie, couleur, mise en page). \\
Message visuel & Contenu transmis par un ensemble d'éléments graphiques organisés. \\
Public cible & Groupe de personnes auquel un message est destiné. \\
Brief & Document ou entretien par lequel un client exprime son besoin, ses objectifs et ses contraintes. \\
Contrainte & Limitation technique, économique, temporelle ou culturelle qui pèse sur un projet. \\
Composition & Organisation des éléments graphiques dans un espace donné. \\
Hiérarchie visuelle & Ordre d'importance des éléments dans une composition. \\
Support & Matériau ou média sur lequel le message est diffusé (papier, écran, objet). \\
Fonction & Rôle que doit jouer un visuel (informer, identifier, promouvoir, etc.). \\
\bottomrule
\end{tabular}
\caption{Vocabulaire essentiel de la section 1.}
\end{table}

\subsection{Exercices de réflexion}

Avant de passer à la section suivante, prends le temps de réfléchir aux questions suivantes. Tu peux noter tes réponses par écrit : cela t'aidera à structurer ta pensée et à te préparer aux exercices pratiques.

\subsubsection{Exercice 1 --- Analyse d'un visuel quotidien}

Choisis un objet graphique que tu croises tous les jours : une affiche publicitaire, un emballage de produit, une couverture de magazine, une interface d'application. Analyse-le en répondant aux questions suivantes :

\begin{enumerate}
    \item Quelle est la \terme{fonction} de ce visuel ? (informer, vendre, identifier, divertir, guider\ldots)
    \item Qui est le \terme{public cible} ? Qu'est-ce qui te le fait penser ?
    \item Quel est le \terme{message principal} ? Peux-tu le résumer en une phrase ?
    \item Quels \terme{éléments graphiques} sont utilisés (typographie, image, couleur, formes) et quel rôle joue chacun ?
    \item Quelles \terme{contraintes} peux-tu deviner (format, support, budget, lisibilité à distance) ?
    \item Ce visuel est-il efficace ? Justifie ta réponse avec au moins deux arguments précis.
\end{enumerate}

\subsubsection{Exercice 2 --- Dessin ou design ?}

Choisis deux images : une illustration que tu trouves belle mais qui ne cherche pas à transmettre un message précis, et une affiche ou une page d'accueil de site web. Pour chacune, écris trois phrases expliquant pourquoi elle relève (ou non) du design graphique en utilisant les quatre piliers : fonction, forme, public, contexte.

\subsubsection{Exercice 3 --- Reformuler la définition}

Sans relire la définition, essaie de la réécrire avec tes propres mots en quelques lignes. Compare ensuite avec la définition générale : as-tu oublié un élément (concevoir, organiser, produire, message, public, éléments graphiques) ?

\begin{erreur}
Confondre « j'aime » et « ça fonctionne ». Un visuel peut te plaire sans remplir sa fonction, et inversement un visuel austère peut être redoutablement efficace. Le client ne te paie pas pour ton goût personnel, mais pour résoudre un problème de communication.
\end{erreur}

\begin{retenir}
\begin{itemize}
    \item Le design graphique \terme{conçoit, organise et produit} des messages visuels pour un public déterminé.
    \item Il se distingue de l'art, du dessin, de la décoration et de la simple maîtrise d'un logiciel.
    \item Quatre piliers guident toute démarche : \terme{fonction, forme, public, contexte}.
    \item Le logiciel est un outil : la valeur du graphiste est dans la conception.
\end{itemize}
\end{retenir}

% ============================================================
\section{Histoire et évolution du design graphique}
% ============================================================

\subsection{Pourquoi connaître l'histoire ?}

Le design graphique n'est pas né avec l'ordinateur. Chaque choix que tu feras (une police, une grille, une affiche à fort contraste) s'inscrit dans une tradition de plusieurs siècles. Connaître cette histoire te permet de comprendre \emph{pourquoi} certaines solutions fonctionnent, d'éviter de réinventer ce qui existe déjà et de nourrir ta culture visuelle (nous y reviendrons au chapitre 5).

Une précision utile : le mot \terme{« graphic design »} est généralement attribué à l'Américain William Addison Dwiggins, qui l'emploie en 1922 pour décrire son propre travail. La pratique, elle, est bien plus ancienne.

\subsection{Des origines à l'imprimerie}

\subsubsection{Les premiers signes}

Dès la préhistoire, l'être humain laisse des traces visuelles sur les parois des grottes. Plus tard, les écritures apparaissent : \terme{hiéroglyphes} égyptiens, \terme{cunéiformes} mésopotamiens, puis les alphabets. Écrire, c'est déjà concevoir un système de signes lisibles par un public. Les manuscrits enluminés du Moyen Âge poussent plus loin la réflexion : mise en page, lettrines, hiérarchie entre texte et images.

\subsubsection{L'invention de l'imprimerie}

Vers 1450, Johannes Gutenberg met au point en Europe l'imprimerie à caractères mobiles. Pour la première fois, un texte peut être reproduit à grande échelle avec une qualité constante. Cette révolution a des conséquences directes sur le métier : apparition des \terme{caractères standardisés}, de la \terme{mise en page} régulière, de la division du travail (graveurs, compositeurs, imprimeurs).

\subsection{Du XIXe siècle aux avant-gardes}

\subsubsection{La révolution industrielle et l'affiche}

Au XIXe siècle, la lithographie en couleur et l'industrialisation de l'imprimerie permettent la naissance de l'\terme{affiche illustrée} et de la publicité moderne. Des artistes comme Jules Chéret, Henri de Toulouse-Lautrec ou Alphonse Mucha produisent des affiches qui envahissent les rues de Paris. Le graphisme devient un phénomène urbain et commercial.

En réaction à la production industrielle, le mouvement \terme{Arts and Crafts} de William Morris défend, à la fin du XIXe siècle, le soin apporté au livre et à l'objet, puis l'\terme{Art nouveau} impose des lignes organiques et une ornementation maîtrisée.

\subsubsection{Le Bauhaus et le modernisme}

En 1919, Walter Gropius fonde le \terme{Bauhaus} à Weimar. L'école défend l'union de l'art, de l'artisanat et de l'industrie, et pose un principe qui nous concerne directement : la forme doit découler de la fonction. Les enseignants et élèves explorent des mises en page asymétriques, des typographies sans empattement et une grande économie de moyens. En 1928, Jan Tschichold publie \anglais{Die neue Typographie}, manifeste d'une typographie moderne, claire et rationnelle. L'école est fermée en 1933 sous la pression du régime nazi, mais son influence se diffuse dans le monde entier.

\subsection{Le XXe siècle : style suisse et identités visuelles}

Après 1945, le \terme{style typographique international}, dit « style suisse », s'impose : grilles rigoureuses, typographie sans empattement, photographie objective, textes alignés à gauche. La police \terme{Helvetica}, dessinée en 1957 par Max Miedinger avec Eduard Hoffmann, en devient le symbole, et Josef Müller-Brockmann en est l'un des maîtres.

Parallèlement, aux États-Unis, des designers comme \terme{Paul Rand} et \terme{Saul Bass} inventent le métier de concepteur d'identités visuelles : logos d'entreprise, génériques de films, systèmes cohérents. Rand travaille notamment pour IBM (dont le logo à bandes apparaît en 1972) ; Bass réalise le célèbre générique du film \anglais{Sueurs froides} (\anglais{Vertigo}, 1958) ainsi que les logos d'AT\&T et de United Airlines.

\begin{defbox}[Étude de cas : le style suisse en pratique]
Une affiche de style suisse repose sur trois choix reconnaissables : une \terme{grille} invisible qui structure la page, une \terme{typographie sans empattement} alignée à gauche, et un \terme{fort contraste d'échelle} entre titre et texte. Ces principes n'ont pas vieilli : on les retrouve aujourd'hui dans de nombreuses interfaces web et identités de marque. Retiens-en la leçon : un bon principe de composition survit aux modes.
\end{defbox}

\subsection{L'ère numérique}

Les années 1980 bouleversent le métier :

\begin{table}[h!]
\centering
\begin{tabular}{@{}p{2.2cm}p{11.8cm}@{}}
\toprule
\textbf{Date} & \textbf{Événement} \\
\midrule
1984 & Lancement du Macintosh d'Apple, premier ordinateur grand public à interface graphique ; naissance du langage PostScript chez Adobe. \\
1985 & PageMaker (Aldus) lance la \terme{PAO} (publication assistée par ordinateur) : le graphiste compose lui-même ses pages. \\
1987 & Première version d'\logiciel{Adobe Illustrator}. \\
1990 & Première version d'\logiciel{Adobe Photoshop}. \\
1991 & Naissance du Web : le design s'étend à l'écran. \\
2010 et après & Généralisation du \anglais{responsive design}, des interfaces mobiles et des systèmes de design. \\
\bottomrule
\end{tabular}
\caption{Quelques repères du design graphique numérique.}
\end{table}

Aujourd'hui, le graphiste travaille sur des supports innombrables : print, web, applications, vidéo, réseaux sociaux, réalité augmentée. Les outils changent, mais les principes (hiérarchie, lisibilité, cohérence) restent les mêmes.

\begin{astuce}
Constitue-toi une frise chronologique personnelle : pour chaque période, colle une image de référence (une affiche, une police, un logo). Cet exercice de 30 minutes t'apprendra plus que dix pages de lecture.
\end{astuce}

\begin{retenir}
\begin{itemize}
    \item Le métier évolue avec les techniques : imprimerie, lithographie, PAO, Web.
    \item Le \terme{Bauhaus} et le \terme{style suisse} fondent une grande partie des principes modernes (fonction, grille, lisibilité).
    \item L'ordinateur a changé les outils, pas les principes.
\end{itemize}
\end{retenir}

% ============================================================
\section{Rôle et importance du design graphique dans la communication}
% ============================================================

\subsection{La communication visuelle est partout}

Réveil, écran de téléphone, emballage du petit-déjeuner, panneaux dans la rue, applications, publicités : en une journée, chacun est exposé à des milliers de messages visuels. Presque tous ont été conçus par quelqu'un. Le design graphique est le plus souvent invisible quand il est réussi : on ne remarque pas la bonne signalétique d'une gare, on remarque l'absence de signalétique.

\subsection{Les fonctions du design graphique}

On distingue classiquement cinq grandes fonctions, qui se combinent souvent dans un même projet :

\begin{table}[h!]
\centering
\begin{tabular}{@{}p{3cm}p{5.2cm}p{5.8cm}@{}}
\toprule
\textbf{Fonction} & \textbf{But} & \textbf{Exemples} \\
\midrule
Informer & Rendre une information claire et accessible. & Infographie, signalétique, plan de métro, rapport annuel. \\
Identifier & Rendre une marque ou une organisation reconnaissable. & Logo, charte graphique, papeterie. \\
Persuader & Convaincre, donner envie, inciter à agir. & Publicité, affiche, page de vente. \\
Guider & Orienter l'utilisateur dans un espace ou une interface. & Panneaux, menus, boutons d'une application. \\
Émouvoir & Créer une ambiance, raconter une histoire. & Affiche de film, couverture de livre, pochette d'album. \\
\bottomrule
\end{tabular}
\caption{Les cinq fonctions du design graphique.}
\end{table}

\subsection{Un enjeu économique et social}

Pour une entreprise, le design graphique est un investissement : il construit la reconnaissance de la marque, la confiance du public et la cohérence entre tous les points de contact (site, emballage, boutique, réseaux sociaux). Pour une institution ou une association, il rend les informations lisibles pour tous. Dans l'espace public, il contribue à la sécurité (signalisation routière, pictogrammes d'évacuation) et à l'accessibilité.

Un exemple historique : pour les Jeux olympiques de Munich en 1972, Otl Aicher conçoit un système de pictogrammes sportifs et de signalétique compréhensible par des visiteurs de toutes langues. Ce travail reste une référence de la communication universelle.

\subsection{Étude de cas : le logo FedEx}

\begin{defbox}[Étude de cas : la flèche cachée de FedEx]
Le logo de FedEx, dessiné en 1994 par Lindon Leader au sein de l'agence Landor, est souvent cité comme un exemple de réussite. Entre le « E » et le « x » se cache une \terme{flèche} formée par l'espace négatif. Elle évoque la vitesse et la précision, sans jamais être imposée au regard. Le logo remplit ainsi plusieurs fonctions à la fois : il identifie, il est très lisible (typographie nette, grands aplats de couleur) et il porte un message (rapidité). Retiens ce principe : une idée simple, intégrée à la forme, vaut plus qu'un effet décoratif.
\end{defbox}

\subsection{Les conséquences d'un design raté}

Un mauvais design n'est pas seulement laid : il a des conséquences concrètes.

\begin{itemize}
    \item Une notice illisible entraîne des erreurs d'utilisation.
    \item Un site mal hiérarchisé fait fuir les visiteurs avant qu'ils trouvent l'information.
    \item Un logo trop complexe devient illisible en petit format.
    \item Un message ambigu peut choquer ou être mal compris.
\end{itemize}

\begin{erreur}
Croire que le design graphique est « la cerise sur le gâteau », à ajouter en fin de projet. Une identité ou une interface pensée trop tard coûte généralement plus cher à corriger qu'à concevoir correctement dès le départ.
\end{erreur}

\begin{retenir}
\begin{itemize}
    \item Le design graphique remplit cinq fonctions : \terme{informer, identifier, persuader, guider, émouvoir}.
    \item Il a un rôle économique (marque), social (accessibilité) et pratique (sécurité, orientation).
    \item Un bon design est souvent invisible ; un mauvais design a des coûts réels.
\end{itemize}
\end{retenir}

% ============================================================
\section{Les différents domaines du design graphique}
% ============================================================

\subsection{Un métier, plusieurs spécialités}

Le design graphique n'est pas un bloc homogène. Selon les supports et les objectifs, il se décline en spécialités qui mobilisent des compétences et des logiciels différents. Beaucoup de graphistes sont généralistes en début de carrière puis se spécialisent ; d'autres restent polyvalents. Connaître la carte des domaines t'aidera à orienter ton apprentissage et, plus tard, ton portfolio.

\subsection{Panorama des domaines}

\begin{table}[h!]
\centering
\begin{tabular}{@{}p{3.3cm}p{6.2cm}p{4.5cm}@{}}
\toprule
\textbf{Domaine} & \textbf{Ce qu'on y conçoit} & \textbf{Outils fréquents} \\
\midrule
Identité visuelle & Logo, charte graphique, papeterie, déclinaisons de marque. & \logiciel{Illustrator} \\
Édition & Livres, magazines, rapports, brochures : mise en page de textes longs. & \logiciel{InDesign} \\
Publicité et affiche & Campagnes, annonces, affiches, bannières. & \logiciel{Photoshop}, \logiciel{Illustrator} \\
Packaging & Emballages, étiquettes, boîtes : design en volume et contraintes d'impression. & \logiciel{Illustrator}, \logiciel{Photoshop} \\
Signalétique & Panneaux, pictogrammes, systèmes d'orientation dans un lieu. & \logiciel{Illustrator} \\
Web et interfaces & Sites, applications, maquettes d'écrans. & \logiciel{Figma}, \logiciel{Photoshop} \\
Motion design & Animations de logos, génériques, vidéos graphiques. & \logiciel{After Effects} \\
Infographie et données & Visualisation d'informations et de statistiques. & \logiciel{Illustrator} \\
Création typographique & Dessin de polices de caractères. & Logiciels de création de fonts \\
\bottomrule
\end{tabular}
\caption{Les principaux domaines du design graphique.}
\end{table}

\subsection{Les domaines en détail}

\subsubsection{L'identité visuelle}

C'est le domaine le plus emblématique : concevoir l'\terme{image d'une organisation}. Il ne s'agit pas seulement d'un logo, mais d'un système : couleurs, typographies, règles d'usage, déclinaisons sur tous les supports. Le livrable principal est la \terme{charte graphique}, que nous étudierons plus loin dans la formation.

\subsubsection{L'édition et la mise en page}

Le graphiste d'édition organise de longs contenus : grilles, styles de paragraphe, hiérarchie des titres, rythme des pages. La qualité se mesure à la \terme{lisibilité} et à la cohérence sur des dizaines ou des centaines de pages.

\subsubsection{La publicité et l'affiche}

Ici, le temps d'attention est très court. Le message doit être compris en quelques secondes : une image forte, une accroche lisible, une signature de marque. C'est le domaine où le lien entre idée créative et exécution graphique est le plus visible.

\subsubsection{Le packaging}

Un emballage doit se démarquer en rayon, informer (composition, mentions légales), protéger le produit et supporter les contraintes de fabrication (découpe, pliage, impression sur des matières variées). Le graphiste travaille à partir d'un \terme{gabarit} fourni par le fabricant.

\subsubsection{Le web et les interfaces}

Le design d'écran ajoute des contraintes propres : tailles d'écran variables, interactions, accessibilité, performance. Les frontières avec l'UX/UI sont proches (voir section 6).

\subsubsection{Le motion design}

Il anime les éléments graphiques dans le temps : transitions, génériques, animations de logo. Il demande en plus une culture du rythme et du montage.

\begin{astuce}
Pour choisir une spécialité, regarde ce que tu aimes \emph{faire} plutôt que ce qui paraît le plus prestigieux. Réalise un mini-projet dans trois domaines différents (un logo, une page de magazine, une affiche) et note celui où le temps passe sans que tu t'en rendes compte.
\end{astuce}

\subsection{Étude de cas : Coca-Cola, une identité qui traverse le temps}

\begin{defbox}[Étude de cas : l'écriture de Coca-Cola]
Le logo manuscrit de Coca-Cola est attribué à Frank Mason Robinson, comptable de l'entreprise, qui l'a tracé en 1886 dans une écriture de type Spencerian. Plus d'un siècle plus tard, cette écriture reste le cœur de l'identité de la marque. Elle montre qu'une identité forte combine plusieurs domaines : logo, typographie, couleur rouge, forme de bouteille, packaging, publicité. La cohérence dans le temps compte autant que l'originalité.
\end{defbox}

\begin{retenir}
\begin{itemize}
    \item Le design graphique se décline en domaines : identité, édition, publicité, packaging, signalétique, web, motion, infographie, typographie.
    \item Chaque domaine a ses contraintes, ses outils et ses livrables.
    \item Une identité forte relie plusieurs domaines dans un même système.
\end{itemize}
\end{retenir}

% ============================================================
\section{Les missions et responsabilités du graphiste}
% ============================================================

\subsection{Ce que fait réellement un graphiste}

Le métier ne se résume pas à « dessiner ». Une journée type mêle écoute du client, recherche, création, itérations, préparation de fichiers, échanges avec les imprimeurs ou les développeurs, et gestion du temps. Les missions suivent le cycle d'un projet.

\subsection{Les missions principales}

\begin{table}[h!]
\centering
\begin{tabular}{@{}p{3.6cm}p{10.4cm}@{}}
\toprule
\textbf{Mission} & \textbf{Description} \\
\midrule
Comprendre le besoin & Analyser le brief, poser les bonnes questions, reformuler la demande. \\
Rechercher & Étudier le secteur, la concurrence, le public, collecter des références. \\
Concevoir & Proposer des pistes, esquisser, définir une direction créative. \\
Réaliser & Produire les maquettes et les éléments graphiques finaux. \\
Présenter et argumenter & Défendre ses choix avec des raisons, pas des goûts personnels. \\
Ajuster & Intégrer les retours du client sans perdre la cohérence du projet. \\
Préparer et livrer & Exporter, vérifier, organiser et transmettre les fichiers adaptés au support. \\
\bottomrule
\end{tabular}
\caption{Les missions d'un graphiste au fil d'un projet.}
\end{table}

\subsection{Les responsabilités}

Le graphiste engage sa responsabilité sur plusieurs plans :

\begin{table}[h!]
\centering
\begin{tabular}{@{}p{3.6cm}p{10.4cm}@{}}
\toprule
\textbf{Domaine} & \textbf{Ce que cela implique} \\
\midrule
Qualité & Fichiers propres, texte relu, résolution et couleurs adaptées au support. \\
Délais et budget & Respecter le planning et le périmètre convenu. \\
Droits d'auteur & N'utiliser que des images, polices et éléments dont la licence autorise l'usage prévu. \\
Éthique & Ne pas tromper le public, ne pas diffuser de contenus discriminatoires. \\
Accessibilité & Garantir une lisibilité pour le plus grand nombre (contraste, taille de texte). \\
Confidentialité & Protéger les informations non publiques du client. \\
\bottomrule
\end{tabular}
\caption{Les responsabilités du graphiste.}
\end{table}

\begin{erreur}
Récupérer une image sur un moteur de recherche et l'utiliser dans un projet client sans vérifier sa licence. La responsabilité du graphiste peut être engagée en cas d'usage non autorisé. Utilise des banques d'images, des polices et des ressources dont la licence permet l'usage commercial, et garde une trace de leurs sources.
\end{erreur}

\subsection{Les différentes situations d'exercice}

\begin{itemize}
    \item \terme{Salarié en agence} : travail en équipe, projets variés, rythme soutenu.
    \item \terme{Salarié en entreprise} (service communication) : cohérence de la marque sur la durée.
    \item \terme{Freelance} : autonomie, mais aussi gestion commerciale, devis et facturation.
    \item \terme{Studio ou collectif} : projets partagés avec d'autres créatifs.
\end{itemize}

Quelle que soit la situation, les mêmes responsabilités s'appliquent.

\begin{retenir}
\begin{itemize}
    \item Les missions suivent le cycle du projet : comprendre, rechercher, concevoir, réaliser, présenter, ajuster, livrer.
    \item Le graphiste est responsable de la qualité, des délais, des droits d'usage, de l'éthique et de l'accessibilité.
    \item Argumenter ses choix est une compétence aussi importante que les créer.
\end{itemize}
\end{retenir}

% ============================================================
\section{Design graphique, illustration, photographie et UI/UX}
% ============================================================

\subsection{Des disciplines voisines, pas identiques}

Le design graphique côtoie d'autres métiers visuels. Les confondre conduit à des malentendus avec les clients et à de mauvais choix de carrière. Les distinguer permet aussi de savoir quand collaborer : un graphiste travaille souvent avec des illustrateurs, des photographes et des designers d'interface.

\subsection{Tableau comparatif}

\begin{table}[h!]
\centering
\begin{tabular}{@{}p{2.6cm}p{3.1cm}p{3.7cm}p{4.6cm}@{}}
\toprule
\textbf{Discipline} & \textbf{Objectif principal} & \textbf{Matière première} & \textbf{Livrable typique} \\
\midrule
Design graphique & Communiquer un message à un public. & Typographie, image, couleur, mise en page. & Logo, affiche, brochure, charte. \\
Illustration & Représenter une idée par le dessin. & Dessin, peinture, techniques numériques. & Image originale pour un livre, un magazine, un emballage. \\
Photographie & Capturer une image du réel ou une scène mise en scène. & Lumière, cadrage, sujet. & Photos, séries, reportages. \\
Design UI/UX & Rendre un produit numérique utilisable et agréable. & Parcours utilisateur, composants d'interface, tests. & Maquettes interactives, prototypes, design systems. \\
\bottomrule
\end{tabular}
\caption{Comparaison de quatre disciplines visuelles.}
\end{table}

\subsection{Les relations entre ces disciplines}

\subsubsection{Design graphique et illustration}

L'illustrateur \emph{produit une image} ; le graphiste \emph{l'intègre dans un ensemble qui communique}. Sur une couverture de livre, l'illustration est un élément, le design graphique est la composition globale (titre, typographie, hiérarchie, format).

\subsubsection{Design graphique et photographie}

Le photographe crée la matière visuelle ; le graphiste la sélectionne, la recadre, la retouche et la met en page. Le retoucheur travaille dans \logiciel{Photoshop} : c'est un point de contact fréquent entre les deux métiers.

\subsubsection{Design graphique et UI/UX}

L'\terme{UX} (expérience utilisateur) s'occupe de la façon dont un produit se comporte et se comprend : parcours, structure, ergonomie. L'\terme{UI} (interface utilisateur) en dessine l'apparence : boutons, couleurs, typographie, icônes. Le design graphique fournit des bases à l'UI (hiérarchie, couleur, typographie), mais l'UX demande des méthodes spécifiques : recherche utilisateur, tests, analyse de données.

\begin{defbox}[À retenir : la question qui tranche]
Pour savoir à quelle discipline appartient une tâche, demande-toi : \emph{« Le but est-il de produire une image, de communiquer un message, ou de faciliter une action ? »} Image : illustration ou photographie. Message : design graphique. Action : UX/UI.
\end{defbox}

\subsection{Les profils hybrides}

Dans la pratique, les frontières sont poreuses. Un graphiste peut illustrer, un illustrateur peut faire de la mise en page, un designer UI peut dessiner des logos. Ce qui compte, c'est de savoir clairement dans quel rôle tu interviens sur un projet donné, et de le dire au client.

\begin{astuce}
Quand un client demande « un site web », demande-lui s'il attend un design graphique (apparence), un design UX (structure et parcours) ou un développement (code). Poser la question dès le départ évite des malentendus sur le périmètre et le prix.
\end{astuce}

\begin{retenir}
\begin{itemize}
    \item Design graphique : \terme{communiquer un message}. Illustration : \terme{créer une image}. Photographie : \terme{capturer ou mettre en scène}. UI/UX : \terme{faciliter l'usage d'un produit numérique}.
    \item Ces disciplines collaborent et se chevauchent, mais n'ont pas le même objectif.
    \item Préciser son rôle sur un projet évite les malentendus.
\end{itemize}
\end{retenir}

% ============================================================
\section{Les qualités et compétences du graphiste professionnel}
% ============================================================

\subsection{Deux familles de compétences}

On distingue les \terme{compétences techniques} (hard skills), qui s'apprennent par la pratique et les formations, et les \terme{compétences humaines} (soft skills), qui font la différence au quotidien avec les clients et les équipes. Un graphiste brillant techniquement mais incapable d'écouter ou de respecter un délai aura du mal à durer.

\subsection{Les compétences techniques}

\begin{table}[h!]
\centering
\begin{tabular}{@{}p{3.6cm}p{10.4cm}@{}}
\toprule
\textbf{Compétence} & \textbf{Ce qu'elle recouvre} \\
\midrule
Composition & Hiérarchie, équilibre, alignement, grille (chapitre 2). \\
Couleur & Harmonies, contraste, modes colorimétriques (chapitre 3). \\
Typographie & Choix, association, lisibilité, hiérarchie des textes (chapitre 4). \\
Culture visuelle & Connaissance des références, capacité d'analyse (chapitre 5). \\
Outils & Maîtrise de \logiciel{Photoshop} et \logiciel{Illustrator}, notions d'\logiciel{InDesign}. \\
Production & Formats, résolution, fonds perdus, export, préparation pour l'impression ou le web. \\
\bottomrule
\end{tabular}
\caption{Les compétences techniques essentielles.}
\end{table}

\subsection{Les qualités humaines}

\begin{table}[h!]
\centering
\begin{tabular}{@{}p{3.6cm}p{10.4cm}@{}}
\toprule
\textbf{Qualité} & \textbf{Pourquoi elle compte} \\
\midrule
Curiosité & Nourrit la créativité et la culture visuelle. \\
Écoute & Comprendre le besoin réel, au-delà de la demande exprimée. \\
Sens critique & Évaluer son propre travail sans complaisance. \\
Rigueur & Éviter les erreurs coûteuses (fautes, fichiers incomplets). \\
Organisation & Tenir les délais, gérer plusieurs projets, nommer et ranger ses fichiers. \\
Communication & Présenter et argumenter ses choix, accepter les retours. \\
Adaptabilité & Les outils et les tendances évoluent en permanence. \\
\bottomrule
\end{tabular}
\caption{Les qualités humaines du graphiste.}
\end{table}

\subsection{Apprendre à progresser}

Aucune de ces compétences ne s'acquiert d'un coup. Voici une démarche simple pour progresser :

\begin{enumerate}
    \item \terme{Pratiquer régulièrement} : de petits exercices fréquents valent mieux qu'un long marathon.
    \item \terme{Analyser} le travail des autres pour en comprendre les choix.
    \item \terme{Demander des retours} et les utiliser.
    \item \terme{Constituer un portfolio} au fur et à mesure.
    \item \terme{Se tenir informé} des évolutions du métier.
\end{enumerate}

\begin{astuce}
Tiens un carnet de références (physique ou numérique) où tu notes, pour chaque visuel qui t'a marqué, \emph{pourquoi} il fonctionne. En quelques mois, tu auras construit un vocabulaire visuel personnel.
\end{astuce}

\begin{erreur}
Croire que la maîtrise de \logiciel{Photoshop} ou \logiciel{Illustrator} suffit à faire un graphiste. Un client peut trouver quelqu'un qui connaît le logiciel ; ce qui a de la valeur, c'est la capacité à comprendre un problème et à y répondre visuellement.
\end{erreur}

\begin{retenir}
\begin{itemize}
    \item Le graphiste combine des compétences \terme{techniques} (composition, couleur, typographie, outils) et \terme{humaines} (écoute, rigueur, communication).
    \item La culture visuelle et la pratique régulière sont les moteurs de la progression.
    \item L'outil ne remplace pas la conception.
\end{itemize}
\end{retenir}

% ============================================================
\section{Le processus de création graphique}
% ============================================================

\subsection{Pourquoi suivre un processus ?}

Un projet réussi n'est jamais le fruit du hasard. Un processus structuré t'évite de partir dans la mauvaise direction, de refaire le travail plusieurs fois et de perdre la confiance du client. Il rassure aussi : le client voit où en est le projet et comprend ce qu'on attend de lui.

\subsection{Les étapes du processus}

\begin{table}[h!]
\centering
\begin{tabular}{@{}p{0.9cm}p{3.3cm}p{9.8cm}@{}}
\toprule
\textbf{N°} & \textbf{Étape} & \textbf{Contenu et livrable} \\
\midrule
1 & Brief & Comprendre la demande : objectifs, public, contraintes, budget, délai. Livrable : brief validé. \\
2 & Recherche & Analyser le secteur, la concurrence, collecter des références. Livrable : moodboard, notes. \\
3 & Idéation & Générer des idées, esquisser rapidement (papier, croquis). Livrable : pistes de concepts. \\
4 & Conception & Développer les meilleures pistes sur ordinateur. Livrable : maquettes. \\
5 & Présentation & Défendre les propositions et recueillir les retours. Livrable : choix de la direction. \\
6 & Ajustements & Affiner selon les retours. Livrable : version validée. \\
7 & Production & Préparer les fichiers finaux, déclinaisons, exports. Livrable : fichiers prêts à l'emploi. \\
8 & Livraison et suivi & Transmettre, vérifier le résultat, archiver. Livrable : dossier de livraison. \\
\bottomrule
\end{tabular}
\caption{Les huit étapes du processus de création.}
\end{table}

\subsection{Les étapes en détail}

\subsubsection{Le brief}

Tout commence par une conversation. Pose des questions simples : à quoi servira ce visuel ? À qui s'adresse-t-il ? Où sera-t-il diffusé ? Quel est le budget et la date limite ? Y a-t-il des éléments imposés (couleurs, logo existant) ? Reformule ensuite la demande par écrit pour vérifier que vous parlez bien de la même chose.

\subsubsection{La recherche et le moodboard}

Un \terme{moodboard} (planche d'ambiance) rassemble images, couleurs, polices et textures qui évoquent l'esprit recherché. C'est un outil de dialogue : il permet de valider une direction avec le client avant de passer des heures sur la conception.

\subsubsection{L'idéation et les croquis}

Avant d'ouvrir un logiciel, dessine. Dix croquis rapides au crayon valent mieux qu'une seule maquette soignée sur écran : tu explores plus de pistes, plus vite, et tu évites de t'attacher à la première idée.

\subsubsection{La conception sur ordinateur}

C'est ici que \logiciel{Illustrator} (vectoriel : logos, pictogrammes) et \logiciel{Photoshop} (retouche et images) interviennent. Travaille avec des calques organisés, des noms clairs et des formats adaptés dès le départ.

\subsubsection{La présentation}

Présente un nombre limité de propositions (généralement deux ou trois), en expliquant \emph{pourquoi} chacune répond au brief. Un argument du type « parce que le public est jeune et sensible à la clarté » est plus convaincant que « parce que je trouve ça joli ».

\subsubsection{La production et la livraison}

Avant de livrer, vérifie : orthographe, résolution, mode colorimétrique (RVB pour l'écran, CMJN pour l'impression), fonds perdus si nécessaire, polices incorporées ou vectorisées selon le cas, nommage des fichiers. Archive le projet pour pouvoir le modifier plus tard.

\begin{defbox}[Étude de cas : le logo Nike]
Le logo « Swoosh » de Nike a été dessiné en 1971 par Carolyn Davidson, alors étudiante en design graphique, à la demande de Phil Knight. Il lui aurait été payé 35 dollars. Cette anecdote est célèbre pour deux raisons : elle montre qu'un symbole très simple peut devenir l'un des plus reconnus au monde, et elle illustre l'importance de \terme{définir la valeur de son travail} dès le départ. Un processus clair (brief, devis, droits cédés) protège à la fois le client et le graphiste.
\end{defbox}

\begin{astuce}
Garde une trace écrite de chaque validation du client (courriel, message). Si un désaccord survient plus tard sur ce qui avait été accepté, tu auras un point de référence.
\end{astuce}

\begin{erreur}
Passer directement à la maquette sur ordinateur sans brief ni recherche. Résultat fréquent : un beau visuel qui ne répond pas à la demande, et plusieurs allers-retours coûteux avec le client.
\end{erreur}

\begin{retenir}
\begin{itemize}
    \item Le processus suit huit étapes : \terme{brief, recherche, idéation, conception, présentation, ajustements, production, livraison}.
    \item Le brief et la recherche préviennent la majorité des erreurs.
    \item Les croquis précèdent l'ordinateur ; les arguments justifient les choix.
    \item Vérifier et archiver avant de livrer fait partie du travail.
\end{itemize}
\end{retenir}

% ============================================================
\section*{Exercices pratiques}
\addcontentsline{toc}{section}{Exercices pratiques}
% ============================================================

Ces trois exercices te demandent d'ouvrir un logiciel pour la première fois dans cette formation. Ne cherche pas la perfection : cherche à appliquer la méthode.

\subsubsection*{Exercice A --- Mini-brief et moodboard (\logiciel{Photoshop})}

\textbf{Consigne :} choisis un projet fictif (une affiche pour un concert local, par exemple). Rédige un brief en cinq lignes (objectif, public, support, contraintes, ton). Dans \logiciel{Photoshop}, crée un document A4 et assemble un moodboard de six à huit images, couleurs et polices qui évoquent l'ambiance voulue.

\textbf{Critère de réussite :} on doit pouvoir comprendre l'ambiance du projet en regardant le moodboard seul, sans lire le brief ; les images sont organisées (alignées, de tailles cohérentes) et sur des calques nommés.

\subsubsection*{Exercice B --- Le logo en trois croquis (\logiciel{Illustrator})}

\textbf{Consigne :} pour une boulangerie imaginaire, dessine au crayon dix croquis rapides de logo, puis sélectionne les trois meilleurs. Reproduis l'un d'eux dans \logiciel{Illustrator} avec les formes de base et l'outil Texte.

\textbf{Critère de réussite :} le logo reste lisible en noir et blanc et réduit à environ 2 cm de large ; il utilise au maximum deux couleurs et une seule famille de police.

\subsubsection*{Exercice C --- Analyse d'un visuel avec les quatre piliers}

\textbf{Consigne :} choisis une affiche publique ou une page d'accueil de site. Dans un document de deux pages (\logiciel{Photoshop} ou \logiciel{Illustrator}), place le visuel et annote-le avec des flèches et des légendes : fonction, forme, public, contexte.

\textbf{Critère de réussite :} chaque pilier est illustré par au moins deux annotations précises et justifiées ; les annotations sont lisibles et ne cachent pas le visuel.

% ============================================================
\section*{Conclusion du chapitre}
\addcontentsline{toc}{section}{Conclusion du chapitre}
% ============================================================

Nous voici au terme du chapitre 1. Tu disposes maintenant des repères essentiels :

\begin{itemize}
    \item \terme{Définition} : le design graphique conçoit, organise et produit des messages visuels pour un public déterminé.
    \item \terme{Histoire} : des manuscrits à l'ère numérique, en passant par l'imprimerie, le Bauhaus et le style suisse.
    \item \terme{Rôle} : informer, identifier, persuader, guider, émouvoir.
    \item \terme{Domaines} : identité, édition, publicité, packaging, signalétique, web, motion, infographie, typographie.
    \item \terme{Missions et responsabilités} : de la compréhension du brief à la livraison, avec rigueur et éthique.
    \item \terme{Disciplines voisines} : illustration, photographie, UI/UX, avec lesquelles on collabore.
    \item \terme{Qualités et compétences} : techniques et humaines.
    \item \terme{Processus} : huit étapes du brief à la livraison.
\end{itemize}

Le design graphique est un métier de méthode autant que de créativité. Les chapitres suivants t'apportent les connaissances de base : composition, couleur, typographie, culture visuelle. Elles te serviront ensuite à utiliser \logiciel{Photoshop} et \logiciel{Illustrator} avec intention plutôt que par imitation.

Dans le chapitre suivant, nous étudierons la \terme{composition graphique} : comment organiser les éléments visuels dans un espace pour produire un message clair, lisible et harmonieux.

\vspace{2em}

\begin{center}
\textit{Fin du chapitre 1 --- Introduction au design graphique.}

\vspace{1em}

\textbf{Prochain chapitre : Chapitre 2 --- La composition graphique.}
\end{center}
